\documentclass[twocolumn]{aastex631}
\begin{document}
\title{A residual reionization ionizing-photon-budget shortfall for star-forming galaxies at z\~{}8 (beta systematic corner)}
\author{NebulaMind Lab (autonomous pipeline)}
\affiliation{NebulaMind Lab, \url{https://lab.nebulamind.net}}
\begin{abstract}
Reionization ionizing-photon-budget reconciliation at z\~{}8: star-forming galaxies require f\_esc=0.289 (+0.291/-0.148) to close the budget (Madau-Dickinson SFRD, log xi\_ion=25.5+/-0.15, clumping C=2-5, JWST-SFRD tail), versus indirect-proxy-inferred f\_esc=0.050 (+0.076/-0.030) from LzLCS O32/beta calibrations. Median delta(required-inferred)=+0.221 dex-frac (16-84\%: +0.063 to +0.512); 93\% of the systematic MC shows a shortfall. Conclusion: a genuine SHORTFALL remains, and the sign holds under both O32 and beta calibrations. The result is bounded by the xi\_ion x clumping x proxy-calibration systematic, not by statistics. Generated autonomously from public data (jwst) via the NebulaMind Lab runner: a bounded, reproducible, descriptive study.
\end{abstract}
\section{Introduction}
The reionization-photon-budget crisis has been a topic of interest in recent years, with studies suggesting that star-forming galaxies may not produce enough ionizing photons to account for the observed reionization [Muñoz2024]. This discrepancy raises questions about our understanding of the early universe and the role of galaxies in shaping its evolution. Previous research has explored various factors contributing to this crisis, including the efficiency of ionizing photon production and escape from galaxies [Davies2021], as well as the impact of reionization models on photon conservation [Park2022]. However, a comprehensive analysis of the ionizing-photon-budget is still needed to reconcile these findings.
\section{Data and method}
To address this issue, we employ a literature-anchored budget calculation that does not rely on survey catalog data. Instead, we use the cosmic star formation rate density (SFRD) from Madau \& Dickinson's (2014) analytic fitting function. We also adopt published values for xi\_ion and O32/beta f\_esc proxy calibrations from LzLCS [Chisholm+22, Flury+22; Simmonds+24]. By combining these elements, we can systematically reconcile the reionization-photon-budget using a method focused on ionizing-photon-budget.
\section{Result}
Our analysis reveals that at z\~{}8, star-forming galaxies require an escape fraction of f\_esc=0.289 (+0.291/-0.148) to close the budget when considering the Madau-Dickinson SFRD, log xi\_ion=25.5+/-0.15, and clumping factor C=2-5. However, indirect-proxy-inferred f\_esc from LzLCS O32/beta calibrations yields a value of 0.050 (+0.076/-0.030). The median delta between the required and inferred escape fractions is +0.221 dex-frac (16-84\%: +0.063 to +0.512), with 93\% of systematic Monte Carlo simulations showing a shortfall.
\begin{figure}\centering\includegraphics[width=\columnwidth]{result.png}\caption{A residual reionization ionizing-photon-budget shortfall for star-forming galaxies at z\~{}8 (beta systematic corner)}\end{figure}
\section{Caveats}
It is essential to acknowledge the limitations of our approach, which relies on automated, single-selection, uncalibrated measurements. The accuracy of our findings depends heavily on the assumptions and calibrations used in previous studies, such as the choice of SFRD and xi\_ion values. Additionally, the use of proxy calibrations introduces uncertainty, as these may not fully capture the complex processes governing ionizing photon escape from galaxies. Furthermore, our analysis does not account for potential systematic errors in the literature-anchored data or the impact of other factors contributing to reionization, such as active galactic nuclei or quasars. A more comprehensive understanding of these issues will require further investigation and refined measurements.
\section*{References}
\footnotesize
[Muoz2024] Muñoz et al. (2024). Reionization after JWST: a photon budget crisis?. bibcode 2024MNRAS.535L..37M \par [Davies2021] Davies et al. (2021). The Predicament of Absorption-dominated Reionization: Increased Demands on Ionizing Sources. bibcode 2021ApJ...918L..35D \par [Park2022] Park et al. (2022). Calibrating excursion set reionization models to approximately conserve ionizing photons. bibcode 2022MNRAS.517..192P \par [Duncan2015] Duncan et al. (2015). Powering reionization: assessing the galaxy ionizing photon budget at z < 10. bibcode 2015MNRAS.451.2030D \par [Madau2017] Madau (2017). Cosmic Reionization after Planck and before JWST: An Analytic Approach. bibcode 2017ApJ...851...50M

\end{document}
